% GENERATED FILE. Edit canonical lessons, then run npm run generate:books.
% canonical-source-hash: fnv1a64:eadf82d0c6d7c19b
% canonical-lessons: KA-C112-samarthya, KA-C112-kausalya, KA-C112-hakku, KA-C112-avakasa, KA-C112-sadhana

\chapter{Ability, Skill, Right, Chance, Tool}
\label{ch:ka-ability-skill-right-chance-tool}
\hlchaptermodality{112}

\begin{chapteropening}
\textbf{By the end of this chapter:} \emph{I can say ability, a skill, a right, a chance and a tool.}

Complete the last lesson of chapter 112: I can say ability, a skill, a right, a chance and a tool.
\end{chapteropening}

\section[\texorpdfstring{\kn{ಸಾಮರ್ಥ್ಯ} (sāmarthya)}{ಸಾಮರ್ಥ್ಯ (sāmarthya)}]{\kn{ಸಾಮರ್ಥ್ಯ} (sāmarthya) — ability}

\label{lesson:KA-C112-samarthya}



\begin{quote}
Before the new one: say the Kannada for a car, then the Kannada for a bicycle.
\end{quote}

\subsection*{You'll want to know: \kn{ಸಾಮರ್ಥ್ಯ}}
\textbf{\kn{ಸಾಮರ್ಥ್ಯ}} — \emph{sāmarthya} — \textquotedblleft{}ability\textquotedblright{}.

\begin{grammarlens}[title={What you've built}]
One more word for this chapter.
\end{grammarlens}

\subsection*{Your turn}
\begin{itemize}
\raggedright
  \item \emph{Say it:} \emph{sāmarthya}
  \item \emph{Say it:} \emph{sāmarthya}, once more
  \item \emph{Say it:} say \emph{sāmarthya}
  \item \emph{Recall:} say the Kannada for a car, then the Kannada for a bicycle, then say \emph{sāmarthya} again
\end{itemize}

\begin{tcolorbox}[breakable,colback=teal!4,colframe=teal!35!black,title={Before you move on}]
What does \kn{ಸಾಮರ್ಥ್ಯ} mean? (\textquotedblleft{}Ability\textquotedblright{}.) Say it once more.
\end{tcolorbox}
\section[\texorpdfstring{\kn{ಕೌಶಲ್ಯ} (kauśalya)}{ಕೌಶಲ್ಯ (kauśalya)}]{\kn{ಕೌಶಲ್ಯ} (kauśalya) — a skill}

\label{lesson:KA-C112-kausalya}



\begin{quote}
Before the new one: say the Kannada for a bicycle, then the Kannada for ability.
\end{quote}

\subsection*{You'll want to know: \kn{ಕೌಶಲ್ಯ}}
\textbf{\kn{ಕೌಶಲ್ಯ}} — \emph{kauśalya} — \textquotedblleft{}a skill\textquotedblright{}.

\begin{grammarlens}[title={What you've built}]
One more word for this chapter.
\end{grammarlens}

\subsection*{Your turn}
\begin{itemize}
\raggedright
  \item \emph{Say it:} \emph{kauśalya}
  \item \emph{Say it:} \emph{kauśalya}, once more
  \item \emph{Say it:} say \emph{kauśalya}
  \item \emph{Recall:} say the Kannada for a bicycle, then the Kannada for ability, then say \emph{kauśalya} again
\end{itemize}

\begin{tcolorbox}[breakable,colback=teal!4,colframe=teal!35!black,title={Before you move on}]
What does \kn{ಕೌಶಲ್ಯ} mean? (\textquotedblleft{}A skill\textquotedblright{}.) Say it once more.
\end{tcolorbox}
\section[\texorpdfstring{\kn{ಹಕ್ಕು} (hakku)}{ಹಕ್ಕು (hakku)}]{\kn{ಹಕ್ಕು} (hakku) — a right}

\label{lesson:KA-C112-hakku}



\begin{quote}
Before the new one: say the Kannada for ability, then the Kannada for a skill.
\end{quote}

\subsection*{You'll want to know: \kn{ಹಕ್ಕು}}
\textbf{\kn{ಹಕ್ಕು}} — \emph{hakku} — \textquotedblleft{}a right\textquotedblright{}.

\begin{grammarlens}[title={What you've built}]
One more word for this chapter.
\end{grammarlens}

\subsection*{Your turn}
\begin{itemize}
\raggedright
  \item \emph{Say it:} \emph{hakku}
  \item \emph{Say it:} \emph{hakku}, once more
  \item \emph{Say it:} say \emph{hakku}
  \item \emph{Recall:} say the Kannada for ability, then the Kannada for a skill, then say \emph{hakku} again
\end{itemize}

\begin{tcolorbox}[breakable,colback=teal!4,colframe=teal!35!black,title={Before you move on}]
What does \kn{ಹಕ್ಕು} mean? (\textquotedblleft{}A right\textquotedblright{}.) Say it once more.
\end{tcolorbox}
\section[\texorpdfstring{\kn{ಅವಕಾಶ} (avakāśa)}{ಅವಕಾಶ (avakāśa)}]{\kn{ಅವಕಾಶ} (avakāśa) — a chance, an opportunity}

\label{lesson:KA-C112-avakasa}



\begin{quote}
Before the new one: say the Kannada for a skill, then the Kannada for a right.
\end{quote}

\subsection*{You'll want to know: \kn{ಅವಕಾಶ}}
\textbf{\kn{ಅವಕಾಶ}} — \emph{avakāśa} — \textquotedblleft{}a chance, an opportunity\textquotedblright{}.

\begin{grammarlens}[title={What you've built}]
One more word for this chapter.
\end{grammarlens}

\subsection*{Your turn}
\begin{itemize}
\raggedright
  \item \emph{Say it:} \emph{avakāśa}
  \item \emph{Say it:} \emph{avakāśa}, once more
  \item \emph{Say it:} say \emph{avakāśa}
  \item \emph{Recall:} say the Kannada for a skill, then the Kannada for a right, then say \emph{avakāśa} again
\end{itemize}

\begin{tcolorbox}[breakable,colback=teal!4,colframe=teal!35!black,title={Before you move on}]
What does \kn{ಅವಕಾಶ} mean? (\textquotedblleft{}A chance, an opportunity\textquotedblright{}.) Say it once more.
\end{tcolorbox}
\section[\texorpdfstring{\kn{ಸಾಧನ} (sādhana)}{ಸಾಧನ (sādhana)}]{\kn{ಸಾಧನ} (sādhana) — a tool, a means}

\label{lesson:KA-C112-sadhana}



\begin{quote}
Before the new one: say the Kannada for a right, then the Kannada for a chance.
\end{quote}

\subsection*{You'll want to know: \kn{ಸಾಧನ}}
\textbf{\kn{ಸಾಧನ}} — \emph{sādhana} — \textquotedblleft{}a tool, a means\textquotedblright{}.

\begin{grammarlens}[title={What you've built}]
One more word for this chapter.
\end{grammarlens}

\subsection*{Your turn}
\begin{itemize}
\raggedright
  \item \emph{Say it:} \emph{sādhana}
  \item \emph{Say it:} \emph{sādhana}, once more
  \item \emph{Say it:} say \emph{sādhana}
  \item \emph{Recall:} say the Kannada for a right, then the Kannada for a chance, then say \emph{sādhana} again
\end{itemize}

\begin{tcolorbox}[breakable,colback=teal!4,colframe=teal!35!black,title={Before you move on}]
What does \kn{ಸಾಧನ} mean? (\textquotedblleft{}A tool, a means\textquotedblright{}.) Say it once more.
\end{tcolorbox}
